% Part IV-B — SMT Integration and Semantic Moves (Slides 81–100)

% ---------------------------------------------------------------------------
% Slide 81: What Is an SMT Solver?
% ---------------------------------------------------------------------------
\begin{frame}
\frametitle{What Is an SMT Solver?}
\begin{center}
\begin{tikzpicture}[
    box/.style={rectangle, rounded corners, draw=jugeoBlue, fill=jugeoLight,
                minimum width=24mm, minimum height=10mm, font=\small\bfseries,
                line width=1pt},
    thy/.style={rectangle, rounded corners, draw=jugeoGreen, fill=jugeoGreen!10,
                minimum width=20mm, minimum height=8mm, font=\scriptsize\bfseries,
                line width=0.8pt},
    arr/.style={-{Stealth[length=3mm]}, thick, jugeoBlue}
  ]
  \node[box, fill=jugeoBlue!15] (sat) at (0,1.5) {SAT Solver};
  \node[thy] (lia) at (-3,0)  {Arithmetic};
  \node[thy] (arr) at (-1,0)  {Arrays};
  \node[thy] (str) at (1,0)   {Strings};
  \node[thy] (bv)  at (3,0)   {Bitvectors};
  \node[box, fill=jugeoBlue!25, minimum width=80mm] (smt) at (0,-1.3) {SMT Solver (e.g.\ Z3 by Microsoft)};
  \draw[arr] (sat) -- (lia);
  \draw[arr] (sat) -- (arr);
  \draw[arr] (sat) -- (str);
  \draw[arr] (sat) -- (bv);
  \node[draw=jugeoBlue, dashed, fit=(sat)(lia)(arr)(str)(bv), inner sep=8pt, rounded corners] {};
\end{tikzpicture}
\end{center}
\vfill
\begin{itemize}
  \item \jugmark{SAT} + \jugmark{theories} = SMT
  \item Automatically \textbf{proves or disproves} logical formulas
  \item Returns \texttt{sat}, \texttt{unsat}, or \texttt{unknown}
\end{itemize}
\end{frame}

% ---------------------------------------------------------------------------
% Slide 82: How F* Uses Z3
% ---------------------------------------------------------------------------
\begin{frame}
\frametitle{How F* Uses Z3}
\begin{center}
\begin{tikzpicture}[
    box/.style={rectangle, rounded corners, draw, minimum width=28mm,
                minimum height=10mm, font=\small\bfseries, line width=1pt},
    arr/.style={-{Stealth[length=3mm]}, thick}
  ]
  \node[box, fill=jugeoOrange!20, draw=jugeoOrange] (fstar) at (-3,0) {F*};
  \node[box, fill=jugeoLight, draw=jugeoBlue]        (blob)  at (0,0)  {One Blob};
  \node[box, fill=jugeoRed!15, draw=jugeoRed]        (z3)    at (3,0)  {Z3};
  \draw[arr, jugeoOrange] (fstar) -- node[above, font=\scriptsize]{VCs} (blob);
  \draw[arr, jugeoBlue]   (blob)  -- (z3);
  \node[below=5mm of z3, font=\small, text=jugeoRed] {sat / unsat / unknown};
\end{tikzpicture}
\end{center}
\vspace{6pt}
\begin{alertblock}{The Problem}
  \begin{itemize}
    \item Sends \textbf{all} verification conditions as a single query
    \item \textcolor{jugeoRed}{No fragment awareness} --- treats linear arithmetic the same as strings
    \item \textcolor{jugeoRed}{No fallback strategy} --- if Z3 says \texttt{unknown}, you're stuck
    \item \textcolor{jugeoRed}{No structured result} --- just a boolean answer
  \end{itemize}
\end{alertblock}
\end{frame}

% ---------------------------------------------------------------------------
% Slide 83: JuGeo's Fragment-Aware Approach
% ---------------------------------------------------------------------------
\begin{frame}
\frametitle{JuGeo's Fragment-Aware Approach}
\begin{center}
\begin{tikzpicture}[
    box/.style={rectangle, rounded corners, draw=jugeoBlue, fill=jugeoLight,
                minimum width=22mm, minimum height=9mm, font=\scriptsize\bfseries,
                line width=1pt},
    arr/.style={-{Stealth[length=3mm]}, thick, jugeoBlue}
  ]
  \node[box, minimum width=30mm, fill=jugeoGreen!15, draw=jugeoGreen] (cls) at (0,1.8) {Classify Formula};
  \node[box] (lia) at (-3.5,0) {QF\_LIA};
  \node[box] (bv)  at (-1.2,0) {QF\_BV};
  \node[box] (str) at (1.2,0)  {STRINGS};
  \node[box] (nl)  at (3.5,0)  {NONLINEAR};
  \node[font=\scriptsize, jugeoBlue] at (0,0) {\ldots};
  \draw[arr] (cls) -- (lia);
  \draw[arr] (cls) -- (bv);
  \draw[arr] (cls) -- (str);
  \draw[arr] (cls) -- (nl);
  \node[box, fill=jugeoGreen!15, draw=jugeoGreen] (t1) at (-3.5,-1.4) {Tactic A};
  \node[box, fill=jugeoGreen!15, draw=jugeoGreen] (t2) at (-1.2,-1.4) {Tactic B};
  \node[box, fill=jugeoGreen!15, draw=jugeoGreen] (t3) at (1.2,-1.4)  {Tactic C};
  \node[box, fill=jugeoGreen!15, draw=jugeoGreen] (t4) at (3.5,-1.4)  {Tactic D};
  \draw[arr] (lia) -- (t1);
  \draw[arr] (bv)  -- (t2);
  \draw[arr] (str) -- (t3);
  \draw[arr] (nl)  -- (t4);
\end{tikzpicture}
\end{center}
\vfill
\begin{exampleblock}{Key Innovation}
  \jugmark{Classify} the formula \textbf{before} sending to Z3.\\
  Different math $\Rightarrow$ different solving strategies.\\
  JuGeo recognizes \jugmark{13 decidable fragments}.
\end{exampleblock}
\end{frame}

% ---------------------------------------------------------------------------
% Slide 84: The 13 SMT Fragments
% ---------------------------------------------------------------------------
\begin{frame}
\frametitle{The 13 SMT Fragments}
\centering
\small
\begin{tabular}{lll}
\toprule
\textbf{Fragment} & \textbf{Domain} & \textbf{Timeout}\\
\midrule
\texttt{QF\_LIA}    & Linear integer arithmetic    & 5\,s\\
\texttt{QF\_LRA}    & Linear real arithmetic       & 5\,s\\
\texttt{QF\_BV}     & Bitvectors                   & 3\,s\\
\texttt{QF\_UF}     & Uninterpreted functions       & 2\,s\\
\texttt{QF\_AUFLIA} & Arrays + UF + LIA            & 8\,s\\
\texttt{QF\_ABV}    & Arrays + bitvectors          & 6\,s\\
\texttt{STRINGS}    & String constraints            & 4\,s\\
\texttt{SEQUENCES}  & Sequence reasoning            & 6\,s\\
\texttt{ARRAYS}     & Pure array theory             & 4\,s\\
\texttt{DATATYPES}  & Algebraic datatypes           & 5\,s\\
\texttt{NONLINEAR}  & Nonlinear arithmetic          & 15\,s\\
\texttt{QUANTIFIED} & Universally quantified        & 20\,s\\
\texttt{MIXED}      & Multi-theory combination      & 10\,s\\
\bottomrule
\end{tabular}
\vspace{4pt}

\textcolor{jugeoGreen}{\small Each fragment gets a tuned Z3 tactic chain and timeout.}
\end{frame}

% ---------------------------------------------------------------------------
% Slide 85: Fragment Classification
% ---------------------------------------------------------------------------
\begin{frame}
\frametitle{Fragment Classification}
\begin{center}
\begin{tikzpicture}[
    stp/.style={rectangle, rounded corners, draw=jugeoBlue, fill=jugeoLight,
                minimum width=30mm, minimum height=9mm, font=\small\bfseries,
                line width=1pt},
    arr/.style={-{Stealth[length=3mm]}, thick, jugeoBlue}
  ]
  \node[stp] (inp) at (0,2.4)  {Formula};
  \node[stp] (s1)  at (-3,0.8) {Inspect Sorts};
  \node[stp] (s2)  at (0,0.8)  {Function Symbols};
  \node[stp] (s3)  at (3,0.8)  {Quantifier Depth};
  \node[stp, fill=jugeoGreen!15, draw=jugeoGreen] (out) at (0,-0.8)
    {Optimal Z3 Tactic Chain};
  \draw[arr] (inp) -- (s1);
  \draw[arr] (inp) -- (s2);
  \draw[arr] (inp) -- (s3);
  \draw[arr] (s1)  -- (out);
  \draw[arr] (s2)  -- (out);
  \draw[arr] (s3)  -- (out);
\end{tikzpicture}
\end{center}
\vfill
\begin{block}{FragmentClassifier}
  Like a \jugmark{doctor diagnosing before prescribing treatment}:\\[4pt]
  Examines sorts, symbols, and quantifiers $\to$ routes to the best strategy.
\end{block}
\end{frame}

% ---------------------------------------------------------------------------
% Slide 86: Five Routing Strategies
% ---------------------------------------------------------------------------
\begin{frame}
\frametitle{Five Routing Strategies}
\begin{center}
\begin{tikzpicture}[
    strat/.style={rectangle, rounded corners=4pt, draw=jugeoBlue, fill=jugeoLight,
                  minimum width=38mm, minimum height=8mm, font=\small, line width=0.8pt,
                  text=jugeoDark},
    every node/.style={inner sep=4pt}
  ]
  \node[strat] (c) at (0, 2.0) {\jugmark{Cheapest} --- minimize cost};
  \node[strat] (f) at (0, 1.0) {\jugmark{Fastest} --- minimize latency};
  \node[strat] (t) at (0, 0.0) {\jugmark{MostTrusted} --- maximize trust};
  \node[strat] (r) at (0,-1.0) {\jugmark{RoundRobin} --- load balance};
  \node[strat, fill=jugeoGreen!15, draw=jugeoGreen]
               (s) at (0,-2.0) {\jugmark{Smart} --- adaptive (recommended)};
\end{tikzpicture}
\end{center}
\vfill
\begin{center}
  \small Choose based on your project's priorities:\\
  production $\to$ \textbf{MostTrusted}\quad
  CI/CD $\to$ \textbf{Fastest}\quad
  exploration $\to$ \textbf{Smart}
\end{center}
\end{frame}

% ---------------------------------------------------------------------------
% Slide 87: Jurisdiction Management
% ---------------------------------------------------------------------------
\begin{frame}
\frametitle{Jurisdiction Management}
\begin{center}
\begin{tikzpicture}[
    zone/.style={rectangle, rounded corners, minimum width=32mm, minimum height=10mm,
                 font=\small\bfseries, line width=1pt, text=white},
    arr/.style={-{Stealth[length=3mm]}, thick, jugeoBlue}
  ]
  \node[zone, draw=jugeoGreen, fill=jugeoGreen] (z3) at (-3,0) {Z3 Decidable};
  \node[zone, draw=jugeoOrange, fill=jugeoOrange] (rt) at (0.7,0) {Runtime Witness};
  \node[zone, draw=jugeoRed, fill=jugeoRed]     (or) at (4.2,0) {Copilot Oracle};
  \draw[arr] (z3) -- node[above, font=\scriptsize, text=jugeoDark]{escalate} (rt);
  \draw[arr] (rt) -- node[above, font=\scriptsize, text=jugeoDark]{escalate} (or);
  \node[below=4mm of z3, font=\scriptsize, text=jugeoGreen]  {Trust 1.0};
  \node[below=4mm of rt, font=\scriptsize, text=jugeoOrange] {Trust 0.7};
  \node[below=4mm of or, font=\scriptsize, text=jugeoRed]    {Trust 0.3};
\end{tikzpicture}
\end{center}
\vfill
\begin{block}{Escalation Chain}
  \begin{enumerate}
    \item Z3 handles decidable fragments $\to$ \textcolor{jugeoGreen}{highest trust}
    \item Outside Z3's jurisdiction? Try \jugmark{runtime witnesses}
    \item Still stuck? Defer to \jugmark{copilot oracle} with lower trust ceiling
  \end{enumerate}
\end{block}
\end{frame}

% ---------------------------------------------------------------------------
% Slide 88: Auditable Proofs from Z3
% ---------------------------------------------------------------------------
\begin{frame}[fragile]
\frametitle{Auditable Proofs from Z3}
\begin{columns}[T]
\begin{column}{0.48\textwidth}
\begin{alertblock}{F* Gets\ldots}
  \begin{itemize}
    \item \texttt{sat}
    \item \texttt{unsat}
    \item \texttt{unknown}
  \end{itemize}
  \centering\small\textcolor{jugeoRed}{That's it.}
\end{alertblock}
\end{column}
\begin{column}{0.48\textwidth}
\begin{exampleblock}{JuGeo Gets\ldots}
  \begin{itemize}
    \item Proof object
    \item Unsat-core extraction
    \item Trust level
    \item Latency metrics
  \end{itemize}
  \centering\small\textcolor{jugeoGreen}{Structured \texttt{Z3Result}}
\end{exampleblock}
\end{column}
\end{columns}
\vfill
\begin{lstlisting}
class Z3Result:
    verdict: Literal["sat","unsat","unknown"]
    proof_object: Optional[Z3Proof]
    unsat_core: list[str]
    trust_level: float        # 0.0 .. 1.0
    latency_ms: float
\end{lstlisting}
\end{frame}

% ---------------------------------------------------------------------------
% Slide 89: What Are Tactics?
% ---------------------------------------------------------------------------
\begin{frame}
\frametitle{What Are Tactics?}
\begin{center}
\begin{tikzpicture}[
    state/.style={rectangle, rounded corners, draw=jugeoBlue, fill=jugeoLight,
                  minimum width=24mm, minimum height=10mm, font=\small,
                  line width=1pt},
    arr/.style={-{Stealth[length=3mm]}, thick, jugeoBlue}
  ]
  \node[state] (g1) at (0,1.5)  {Goal State A};
  \node[state] (g2) at (-2,0)   {Sub-goal B};
  \node[state] (g3) at (2,0)    {Sub-goal C};
  \node[state, fill=jugeoGreen!15, draw=jugeoGreen] (qed) at (0,-1.5) {QED};
  \draw[arr] (g1) -- node[left, font=\scriptsize]{split} (g2);
  \draw[arr] (g1) -- node[right, font=\scriptsize]{split} (g3);
  \draw[arr] (g2) -- node[left, font=\scriptsize]{simp} (qed);
  \draw[arr] (g3) -- node[right, font=\scriptsize]{auto} (qed);
\end{tikzpicture}
\end{center}
\vfill
\begin{block}{Definition}
  \jugmark{Tactics} are commands that manipulate proof state:
  \begin{itemize}
    \item ``Try this strategy,'' ``split into cases,'' ``apply this lemma''
    \item They are proof \jugmark{scripts} --- sequences of instructions
  \end{itemize}
\end{block}
\end{frame}

% ---------------------------------------------------------------------------
% Slide 90: Tactics in LEAN / Coq / F*
% ---------------------------------------------------------------------------
\begin{frame}
\frametitle{Tactics in LEAN / Coq / F*}
\begin{center}
\small
\begin{tabular}{lll}
\toprule
\textbf{System} & \textbf{Mechanism} & \textbf{Example Tactics}\\
\midrule
LEAN   & \texttt{tactic} mode    & \texttt{unfold, induction, simp, omega}\\
Coq    & Ltac / Ltac2            & \texttt{intro, apply, rewrite, auto}\\
F*     & Meta-F* \texttt{Tac}    & \texttt{norm, smt, trefl, dump}\\
\bottomrule
\end{tabular}
\end{center}
\vspace{8pt}
\begin{alertblock}{Limitation: They Only Touch Proof Terms}
  \begin{itemize}
    \item Cannot manipulate \textcolor{jugeoRed}{\textbf{trust levels}}
    \item Cannot manipulate \textcolor{jugeoRed}{\textbf{evidence quality}}
    \item Cannot trigger \textcolor{jugeoRed}{\textbf{repair}} when proofs break
    \item No notion of \textcolor{jugeoRed}{\textbf{geometric structure}}
  \end{itemize}
\end{alertblock}
\end{frame}

% ---------------------------------------------------------------------------
% Slide 91: JuGeo's Semantic Moves
% ---------------------------------------------------------------------------
\begin{frame}[fragile]
\frametitle{JuGeo's Semantic Moves}
\begin{exampleblock}{Replace Tactics with Semantic Moves}
  First-class proof operations on the \jugmark{full judgment geometry}.
\end{exampleblock}
\vspace{4pt}
\begin{lstlisting}
@dataclass
class SemanticMove:
    id: str
    kind: MoveKind
    target: Coordinate
    preconditions: list[Predicate]
    postconditions: list[Predicate]
    expected_gain: float
    cost_estimate: float
    trust_floor: float         # minimum trust required
\end{lstlisting}
\vfill
\begin{center}
\small
Moves operate on \jugmark{geometry + trust + evidence},\\
not just proof terms.
\end{center}
\end{frame}

% ---------------------------------------------------------------------------
% Slide 92: Move Vocabulary — Structural Rules
% ---------------------------------------------------------------------------
\begin{frame}
\frametitle{Move Vocabulary --- Structural Rules}
\begin{center}
\begin{tikzpicture}[
    mv/.style={rectangle, rounded corners=3pt, draw=jugeoBlue, fill=jugeoLight,
               minimum width=28mm, minimum height=9mm, font=\small\bfseries,
               line width=0.8pt},
    desc/.style={font=\scriptsize, text=jugeoDark, text width=40mm, align=left}
  ]
  \node[mv] (w) at (-3, 1.2) {WEAKEN};
  \node[mv] (c) at (-3,-0.2) {CONTRACT};
  \node[mv] (e) at (-3,-1.6) {EXCHANGE};
  \node[mv, fill=jugeoOrange!15, draw=jugeoOrange] (cu) at (-3,-3.0) {CUT};
  \node[desc, right=6mm of w]  {Add an unused hypothesis};
  \node[desc, right=6mm of c]  {Merge duplicate hypotheses};
  \node[desc, right=6mm of e]  {Reorder hypotheses};
  \node[desc, right=6mm of cu] {Introduce intermediate lemma};
\end{tikzpicture}
\end{center}
\vfill
\begin{center}
\small These correspond to \jugmark{classical structural rules} of sequent calculus.\\
CUT is the most powerful --- and the most dangerous (see Slide 100).
\end{center}
\end{frame}

% ---------------------------------------------------------------------------
% Slide 93: Move Vocabulary — Logical Rules
% ---------------------------------------------------------------------------
\begin{frame}
\frametitle{Move Vocabulary --- Logical Rules}
\begin{center}
\begin{tikzpicture}[
    mv/.style={rectangle, rounded corners=3pt, draw=jugeoBlue, fill=jugeoLight,
               minimum width=30mm, minimum height=9mm, font=\small\bfseries,
               line width=0.8pt},
    desc/.style={font=\scriptsize, text=jugeoDark, text width=42mm, align=left}
  ]
  \node[mv] (v) at (-3, 1.8) {VERIFY};
  \node[mv] (c) at (-3, 0.6) {CONSTRUCT};
  \node[mv] (n) at (-3,-0.6) {NORMALIZE};
  \node[mv] (i) at (-3,-1.8) {INTRODUCE};
  \node[mv] (e) at (-3,-3.0) {ELIMINATE};
  \node[desc, right=6mm of v] {Send sub-goal to Z3};
  \node[desc, right=6mm of c] {Synthesize a witness / solution};
  \node[desc, right=6mm of n] {Simplify; strip irrelevant diffs};
  \node[desc, right=6mm of i] {Introduce connective (e.g.\ $\forall$-intro)};
  \node[desc, right=6mm of e] {Eliminate connective (e.g.\ $\exists$-elim)};
\end{tikzpicture}
\end{center}
\vfill
\begin{center}
\small
VERIFY bridges \jugmark{semantic moves} and the \jugmark{SMT layer}:\\
it dispatches the classified fragment to Z3.
\end{center}
\end{frame}

% ---------------------------------------------------------------------------
% Slide 94: Move Vocabulary — Geometric Moves (NEW!)
% ---------------------------------------------------------------------------
\begin{frame}
\frametitle{Move Vocabulary --- Geometric Moves {\small\textcolor{jugeoGreen}{(NEW!)}}}
\begin{center}
\begin{tikzpicture}[
    mv/.style={rectangle, rounded corners=3pt, draw=jugeoGreen, fill=jugeoGreen!12,
               minimum width=30mm, minimum height=9mm, font=\small\bfseries,
               line width=1pt},
    desc/.style={font=\scriptsize, text=jugeoDark, text width=44mm, align=left}
  ]
  \node[mv] (r) at (-3, 1.5) {RESTRICT};
  \node[mv] (t) at (-3, 0.3) {TRANSPORT};
  \node[mv] (g) at (-3,-0.9) {GLUE};
  \node[mv] (c) at (-3,-2.1) {REFINE\_COVER};
  \node[desc, right=6mm of r] {Narrow to a sub-coordinate};
  \node[desc, right=6mm of t] {Move evidence along a morphism};
  \node[desc, right=6mm of g] {Stitch local sections into global proof};
  \node[desc, right=6mm of c] {Improve the covering of a space};
\end{tikzpicture}
\end{center}
\vfill
\begin{alertblock}{No Analogue in LEAN / Coq / F*}
  These moves exist because JuGeo has \jugmark{geometric structure}:\\
  coordinates, covers, restriction maps, and gluing conditions.
\end{alertblock}
\end{frame}

% ---------------------------------------------------------------------------
% Slide 95: Move Vocabulary — Evidence Moves (NEW!)
% ---------------------------------------------------------------------------
\begin{frame}
\frametitle{Move Vocabulary --- Evidence Moves {\small\textcolor{jugeoGreen}{(NEW!)}}}
\begin{center}
\begin{tikzpicture}[
    mv/.style={rectangle, rounded corners=3pt, draw=jugeoGreen, fill=jugeoGreen!12,
               minimum width=34mm, minimum height=9mm, font=\small\bfseries,
               line width=1pt},
    desc/.style={font=\scriptsize, text=jugeoDark, text width=42mm, align=left}
  ]
  \node[mv] (p) at (-3, 1.2) {PROMOTE\_TRUST};
  \node[mv] (c) at (-3, 0.0) {CHALLENGE};
  \node[mv] (e) at (-3,-1.2) {ESCALATE};
  \node[desc, right=6mm of p] {Upgrade trust with justification};
  \node[desc, right=6mm of c] {Question existing evidence};
  \node[desc, right=6mm of e] {Defer to oracle (lower trust)};
\end{tikzpicture}
\end{center}
\vfill
\begin{exampleblock}{Operating on Trust, Not Just Terms}
  \begin{itemize}
    \item Traditional tactics: transform proof \textbf{terms}
    \item Evidence moves: transform \jugmark{trust levels} and \jugmark{evidence quality}
    \item CHALLENGE can \textit{downgrade} trust if evidence is suspect
  \end{itemize}
\end{exampleblock}
\end{frame}

% ---------------------------------------------------------------------------
% Slide 96: Move Vocabulary — Treaty Moves (NEW!)
% ---------------------------------------------------------------------------
\begin{frame}
\frametitle{Move Vocabulary --- Treaty Moves {\small\textcolor{jugeoGreen}{(NEW!)}}}
\begin{center}
\begin{tikzpicture}[
    mv/.style={rectangle, rounded corners=3pt, draw=jugeoGreen, fill=jugeoGreen!12,
               minimum width=40mm, minimum height=9mm, font=\small\bfseries,
               line width=1pt},
    desc/.style={font=\scriptsize, text=jugeoDark, text width=40mm, align=left}
  ]
  \node[mv] (n) at (-3, 0.6) {NEGOTIATE\_TREATY};
  \node[mv] (r) at (-3,-0.8) {REPAIR};
  \node[desc, right=6mm of n] {Reconcile conflicts between overlapping code sections};
  \node[desc, right=6mm of r] {Fix obstructions when gluing fails};
\end{tikzpicture}
\end{center}
\vfill
\begin{block}{When Code Sections Overlap}
  \begin{itemize}
    \item Module A and Module B both define \texttt{validate()}
    \item They \jugmark{disagree} on return type
    \item \textbf{NEGOTIATE\_TREATY}: find common ground
    \item \textbf{REPAIR}: patch the obstruction so gluing succeeds
  \end{itemize}
\end{block}
\end{frame}

% ---------------------------------------------------------------------------
% Slide 97: 9 Moves With No Analogue
% ---------------------------------------------------------------------------
\begin{frame}
\frametitle{9 Moves With No Analogue in LEAN / Coq / F*}
\begin{center}
\begin{tikzpicture}[
    cat/.style={rectangle, rounded corners=3pt, minimum width=26mm, minimum height=7mm,
                font=\scriptsize\bfseries, line width=0.8pt, text=white},
  ]
  % Geometric
  \node[cat, fill=jugeoBlue] (g1) at (-3.5, 1.8) {RESTRICT};
  \node[cat, fill=jugeoBlue] (g2) at (-0.5, 1.8) {TRANSPORT};
  \node[cat, fill=jugeoBlue] (g3) at (2.5, 1.8)  {GLUE};
  \node[cat, fill=jugeoBlue] (g4) at (5.2, 1.8)  {REFINE\_COVER};
  \node[above=2mm of g2, font=\small\bfseries, text=jugeoBlue] {Geometric};

  % Evidence
  \node[cat, fill=jugeoGreen] (e1) at (-2.5, 0.2) {PROMOTE\_TRUST};
  \node[cat, fill=jugeoGreen] (e2) at (0.8, 0.2)  {CHALLENGE};
  \node[cat, fill=jugeoGreen] (e3) at (3.8, 0.2)  {ESCALATE};
  \node[above=2mm of e2, font=\small\bfseries, text=jugeoGreen] {Evidence};

  % Treaty
  \node[cat, fill=jugeoOrange] (t1) at (-1.5,-1.3) {NEGOTIATE\_TREATY};
  \node[cat, fill=jugeoOrange] (t2) at (2.3,-1.3)  {REPAIR};
  \node[above=2mm of t1, font=\small\bfseries, text=jugeoOrange, xshift=14mm] {Treaty};
\end{tikzpicture}
\end{center}
\vfill
\begin{center}
\small
These \jugmark{9 moves} exist because JuGeo has geometric structure,\\
graded trust, and inter-module treaty management\\
--- concepts absent from traditional proof assistants.
\end{center}
\end{frame}

% ---------------------------------------------------------------------------
% Slide 98: Adaptive Move Selection
% ---------------------------------------------------------------------------
\begin{frame}
\frametitle{Adaptive Move Selection}
\begin{center}
\begin{tikzpicture}[
    stp/.style={rectangle, rounded corners, draw=jugeoBlue, fill=jugeoLight,
                minimum height=8mm, font=\scriptsize\bfseries, line width=0.8pt,
                text width=26mm, align=center},
    arr/.style={-{Stealth[length=2.5mm]}, thick, jugeoBlue}
  ]
  \node[stp] (en) at (0, 2.6)  {MoveEnumerator};
  \node[stp] (pc) at (0, 1.6)  {PreconditionChecker};
  \node[stp] (pr) at (0, 0.6)  {MovePrioritizer};
  \node[stp] (cr) at (0,-0.4)  {ConflictResolver};
  \node[stp] (ae) at (0,-1.4)  {ApplicationEngine};
  \node[stp, fill=jugeoGreen!15, draw=jugeoGreen] (pv) at (0,-2.4)
    {PostconditionVerifier};
  \draw[arr] (en) -- (pc);
  \draw[arr] (pc) -- (pr);
  \draw[arr] (pr) -- (cr);
  \draw[arr] (cr) -- (ae);
  \draw[arr] (ae) -- (pv);
  % feedback
  \draw[arr, jugeoOrange, dashed] (pv.east) -- ++(1.5,0) |- (en.east)
    node[pos=0.25, right, font=\scriptsize, text=jugeoOrange]{feedback};
\end{tikzpicture}
\end{center}
\vfill
\begin{center}
\small
Not a fixed script --- an \jugmark{adaptive control loop}.\\
Failed postconditions feed back to enumeration.
\end{center}
\end{frame}

% ---------------------------------------------------------------------------
% Slide 99: Control Strategies
% ---------------------------------------------------------------------------
\begin{frame}
\frametitle{Control Strategies}
\begin{center}
\begin{tikzpicture}[
    strat/.style={rectangle, rounded corners=3pt, draw=jugeoBlue, fill=jugeoLight,
                  minimum width=28mm, minimum height=8mm, font=\small, line width=0.8pt,
                  text=jugeoDark},
    desc/.style={font=\scriptsize, text=jugeoDark, text width=45mm, align=left}
  ]
  \node[strat] (g) at (-3, 1.5) {\jugmark{GREEDY}};
  \node[strat] (l) at (-3, 0.3) {\jugmark{LOOKAHEAD}};
  \node[strat] (b) at (-3,-0.9) {\jugmark{BALANCED}};
  \node[strat, fill=jugeoGreen!15, draw=jugeoGreen]
               (a) at (-3,-2.1) {\jugmark{ADAPTIVE}};
  \node[desc, right=6mm of g] {Always pick highest expected gain};
  \node[desc, right=6mm of l] {Simulate $k$ steps ahead, then choose};
  \node[desc, right=6mm of b] {Trade off gain vs.\ cost};
  \node[desc, right=6mm of a] {Switch strategy mid-proof based on progress};
\end{tikzpicture}
\end{center}
\vfill
\begin{exampleblock}{More Powerful Than Meta-F* Combinators}
  \begin{itemize}
    \item F*: \texttt{or\_else}, \texttt{seq}, \texttt{repeat} --- combinators on tactics
    \item JuGeo: full \jugmark{strategy switching} with lookahead and cost models
  \end{itemize}
\end{exampleblock}
\end{frame}

% ---------------------------------------------------------------------------
% Slide 100: Cut Elimination — A Key Theorem
% ---------------------------------------------------------------------------
\begin{frame}
\frametitle{Cut Elimination --- A Key Theorem}
\begin{center}
\begin{tikzpicture}[
    pf/.style={rectangle, rounded corners, draw=jugeoBlue, fill=jugeoLight,
               minimum width=28mm, minimum height=10mm, font=\small\bfseries,
               line width=1pt},
    arr/.style={-{Stealth[length=3mm]}, ultra thick, jugeoGreen}
  ]
  \node[pf] (with) at (-2.5,0) {Proof with CUT};
  \node[pf, fill=jugeoGreen!15, draw=jugeoGreen] (wo) at (2.5,0)
    {Cut-free Proof};
  \draw[arr] (with) -- node[above, font=\scriptsize, text=jugeoDark]
    {normalize} (wo);
\end{tikzpicture}
\end{center}
\vspace{6pt}
\begin{block}{Theorem (Cut Admissibility)}
  Every proof in JuGeo that uses \jugmark{CUT} can be transformed into\\
  an equivalent proof that does \textbf{not} use CUT.
\end{block}
\vspace{4pt}
\begin{exampleblock}{Why This Matters}
  \begin{itemize}
    \item \jugmark{Consistency}: no circular reasoning smuggled via cuts
    \item \jugmark{Normalization}: every proof has a canonical form
    \item \jugmark{Sub-formula property}: conclusions only involve sub-formulas of hypotheses
    \item This is JuGeo's analogue of strong normalization in type theory
  \end{itemize}
\end{exampleblock}
\end{frame}
